\chapter{まとめと今後の課題}

\section{まとめ}

本研究では，
髭の存在が皮膚表面下散乱に与える影響を，
実用的な計算コストで表現することを目的として，
髭密度を用いた表面下散乱モデルを提案した。

従来のアプローチでは，
髭を一本一本の幾何構造として扱う必要があり，
その結果，
計算コストの増大や，
配置条件に強く依存した不安定な挙動が問題となる。
これに対し，
本研究では，
髭の分布を表面上のスカラー量である「髭密度」として近似し，
皮膚モデルと髭モデルを線形に混合することで，
表面下散乱特性を制御する手法を採用した。

評価実験では，
点光源入射時における放射輝度分布 $R_d(r)$ を用いて，
髭なしの場合，
幾何的な髭モデルを含む場合，
および提案した髭密度モデルの場合を比較した。
その結果，
髭の存在が表面下散乱の空間的な広がりに影響を与えることが確認され，
特に髭密度モデルにおいては，
分布形状が不自然に破綻することなく，
髭密度の変化に応じて連続的に制御できることが示された。

また，
幾何的な髭モデルとの比較から，
提案モデルは髭の詳細な幾何構造を厳密に再現するものではないものの，
髭の存在による表面下散乱特性の変化を，
方向性として適切に捉えていることが確認された。
このことから，
髭密度という抽象化は，
表面下散乱に対する髭の影響を安定かつ実用的に扱うための
有効な近似であると結論付けられる。

以上より，
本研究は，
髭のような微細構造を持つ要素を含む皮膚レンダリングにおいて，
密度に基づく近似モデルが有効であることを示し，
実用的な表面下散乱表現に対する一つの設計指針を提示した。

\section{今後の課題}

本研究で提案した髭密度モデルは，
計算コストと表現力のバランスを重視した近似モデルであり，
いくつかの課題が残されている。

まず，
本モデルでは，
髭の影響をスカラーの密度値として表現しているため，
髭の方向性や太さのばらつき，
髭同士の相互遮蔽といった効果は明示的に扱っていない。
これらの要素を考慮することで，
より多様な髭表現に対応できる可能性がある。

次に，
幾何的な髭モデルとの対応関係については，
本研究では主に定性的および差分的な比較に留まっている。
将来的には，
髭の太さや埋め込み深さを系統的に変化させた幾何モデルとの比較を行い，
幾何的要因と髭密度との関係をより定量的に分析することが課題として挙げられる。

また，
本研究では静的な髭密度マップを用いたが，
実際の顔表現では，
部位ごとの密度差や，
成長・剃毛・スタイリングによる時間的変化が存在する。
提案モデルは，
密度マップを差し替えることでこれらの変化に対応可能であるため，
動的な髭密度を扱う拡張も今後の課題である。

さらに，
本手法は髭を対象として検討を行ったが，
同様の考え方は，
産毛や皮膚表面の微細な繊維構造など，
他の微細要素を含む表現にも応用できる可能性がある。
これらへの適用可能性を検討することも，
今後の展開として考えられる。

以上の課題に取り組むことで，
髭密度に基づく表面下散乱モデルの表現力と汎用性をさらに高めることが期待される。
